\section{Computational Complexity and Quantum Advantage}
To really understand what makes a quantum computer computationally "powerful," we need to define the mathematical boundaries of what it can actually do. 
Instead of relying on vague claims, we place quantum algorithms into the well-established framework of classical computational complexity theory \cite{arora2009computational}. 
This gives us an objective, rigorous way to compare how quantum algorithms scale against the absolute best classical methods as the size of a problem grows.

\subsection{Definitions of Complexity Classes}
Before we can map out the territory of quantum advantage, we need to define our key complexity classes. 
Think of these classes as buckets that group computational problems based on how quickly their time and memory requirements blow up as the size of the input, $N$, gets larger.

\begin{itemize}
    \item \textbf{P (Polynomial Time):} 
    The bread and butter of classical computing. 
    This is the class of problems that a standard computer can solve efficiently, meaning the runtime scales gracefully (like $N^2$ or $N^3$).

    \item \textbf{ZPP (Zero-error Probabilistic Polynomial time):} 
    Imagine a classical computer with a random coin. 
    It uses randomness to speed things up, and when it finally gives an answer, it is always 100\% correct. 
    However, we only guarantee that it finishes quickly \textit{on average}. 
    Formally, $\text{ZPP} = \text{RP} \cap \text{co-RP}$.

    \item \textbf{RP (Randomized Polynomial time):} 
    A classical algorithm that runs quickly and uses randomness, but is slightly pessimistic. 
    If the true answer is "no", it always correctly says "no". 
    If the true answer is "yes", it might accidentally say "no", but it has at least a 50\% chance of getting the "yes" right.

    \item \textbf{BPP (Bounded-error Probabilistic Polynomial time):} 
    This is the gold standard for practical classical algorithms. 
    The computer runs quickly but might make a mistake on both "yes" and "no" answers. 
    However, the probability of being right is always bounded strictly away from a coin flip (typically defined as $\geq 2/3$). 
    Because it's reliably better than 50/50, we can just run the program a few times and take a majority vote to confidently find the truth \cite{arora2009computational}.

    \item \textbf{NP (Nondeterministic Polynomial time):} 
    These are problems that might be incredibly hard to solve from scratch, but if someone hands you the correct answer (a "witness"), you can verify it very quickly. 
    Think of a Sudoku puzzle: hard to fill out, but trivial to check if it's right.

    \item \textbf{BQP (Bounded-error Quantum Polynomial time):} 
    The quantum sibling to BPP. 
    This is the exact class of problems that a quantum computer can solve efficiently. 
    It allows for a small margin of error (at most $1/3$), but just like BPP, we can easily shrink that error to near zero by running the quantum circuit a few extra times and taking a vote \cite{bernstein1997, nielsen2010quantum}.

    \item \textbf{PP (Probabilistic Polynomial time):} 
    A bizarre and incredibly powerful theoretical class. 
    A classical computer guesses the answer, and it only needs to be right strictly more than 50\% of the time—even if it's 50.0000001\%. 
    Because the margin is so razor-thin, taking a majority vote would physically require an exponential number of runs. 
    Despite being physically impractical, it is a mathematical behemoth that easily swallows up both NP and BQP \cite{adleman1997}.

    \item \textbf{PSPACE (Polynomial Space):} 
    Problems that can be solved using a reasonable (polynomial) amount of memory. 
    The catch? There is no time limit, meaning the computer might take until the heat death of the universe to finish its calculation.
\end{itemize}

\begin{figure}[hbt!]
    \centering
    \includegraphics[width=0.8\textwidth]{Randomised_Complexity_Classes.jpg}
    \caption{Euler diagram illustrating the theoretical relationships and hierarchy of classical and quantum computational complexity classes. Source: Wikipedia}
    \label{fig:complexity_euler}
\end{figure}

\subsection{Proofs of Inclusion}
Now that we have our classes, we can conceptually stack them like Russian nesting dolls: $\text{P} \subseteq \text{ZPP} \subseteq \text{RP} \subseteq \text{BPP} \subseteq \text{BQP} \subseteq \text{PP} \subseteq \text{PSPACE}$. 
While almost every computer scientist firmly believes that BQP is strictly mathematically larger than BPP—implying quantum computers can efficiently solve problems that classical computers inherently cannot—proving this strict separation remains one of the greatest open problems in theoretical computer science. 
We can, however, definitively prove the inclusion boundaries. 
Let's walk through the intuition for why these theoretical dolls fit inside each other.

\textbf{Theorem 1:} $\text{P} \subseteq \text{BPP}$
\begin{proof}
This one is almost trivially true. 
Suppose we have a problem that can be solved perfectly by a standard deterministic computer (Class P). 
We can simply hand this exact same problem to a randomized computer (Class BPP) and instruct it to never roll its dice. 
Because the program never uses randomness, it behaves exactly like the deterministic machine, succeeding with 100\% certainty. 
Since 100\% is well above the 66\% success threshold required by BPP, the problem comfortably fits inside it.
\end{proof}

\textbf{Theorem 2:} $\text{BPP} \subseteq \text{BQP}$
\begin{proof}
Can a quantum computer simulate a randomized classical one? Yes, and the process requires exactly two conceptual steps.

First, we must deal with the fact that quantum mechanics is fundamentally reversible. 
Standard classical computers throw away information (for example, an AND gate irreversibly crushes two bits into one). 
However, Bennett's theorem mathematically proves that any classical program can be rewritten into a perfectly reversible one using energy-conserving logic like the Toffoli gate, incurring only a tiny structural overhead \cite{bennett1973}. 

Second, the classical BPP machine needs random coin flips to operate. 
A quantum computer is arguably nature's perfect random number generator. 
By taking a fresh qubit in the $|0\rangle$ state and hitting it with a Hadamard gate, we create a perfect 50/50 superposition:
\begin{equation}
    H|0\rangle = \frac{1}{\sqrt{2}}(|0\rangle + |1\rangle)
\end{equation}
Measuring this qubit yields a perfectly unbiased, fundamentally unpredictable random bit. 

By feeding a steady stream of these quantum coin flips into our reversible classical circuit, the quantum computer flawlessly mimics the classical BPP machine. 
It preserves the exact same probability of success, definitively proving that a quantum computer can do everything a randomized classical computer can do.
\end{proof}

\textbf{Theorem 3:} $\text{BQP} \subseteq \text{PSPACE}$
\begin{proof}
This theorem asks a fascinating question: could a classical computer simulate a quantum one if we gave it unlimited time, but strictly limited memory? 
The answer is yes, using an intuitive method inspired by Richard Feynman's sum-over-histories \cite{bernstein1997}.

\begin{figure}[hbt!]
    \centering
    \includegraphics[width=0.7\textwidth]{Sum_of_histories_tree.png}
    \caption{A computational tree representing the Feynman sum-over-histories approach for evaluating the probability amplitudes of quantum circuits. Source: Wikipedia}
    \label{fig:sum_of_histories}
\end{figure}

Imagine the execution of a quantum circuit as a massive branching tree of possibilities. 
If a circuit is made of $m$ sequential gates $U_1, U_2, \dots, U_m$, the transition amplitude from a starting state $|x\rangle$ to a final state $|y\rangle$ is mathematically defined as:
\begin{equation}
    \langle y | U_m U_{m-1} \dots U_1 | x \rangle
\end{equation}

By breaking this down, we can view it as a sum over every possible intermediate state, $z$, at every step of the computation:
\begin{equation}
    \sum_{z_1, \dots, z_{m-1}} \langle y | U_m | z_{m-1} \rangle \langle z_{m-1} | U_{m-1} | z_{m-2} \rangle \dots \langle z_1 | U_1 | x \rangle
\end{equation}

At every single gate operation, the universe of the simulation splits into different paths, each carrying a tiny piece of this complex probability amplitude. 
Because a quantum register with $n$ qubits evaluating $m$ gates creates $2^{nm}$ possible branches, a classical computer would take practically forever to calculate this massive sum. 

However, it does not need to store that entire branching universe in its RAM at once. 
Instead, it can use a simple depth-first search. 
The classical computer traces one single, isolated path all the way from the root to the leaf, calculates its specific amplitude, adds it to a running numerical tally, and then completely overwrites that path in memory to calculate the next one. 

To pull this off, the classical computer only needs to remember two things: the chronological sequence of the current path, and the running probability total. 
Because the circuit size is polynomial, both of these metrics require very few bits to track—scaling only polynomially with the size of the input. 
Since the active memory ceiling stays low even as the execution time stretches toward infinity, the entire simulation comfortably operates within PSPACE.
\end{proof}

\textbf{Theorem 4:} $\text{BQP} \subseteq \text{PP}$
\begin{proof}
As we just saw, calculating a quantum probability requires adding up a staggering number of paths. 
Some of these paths have positive phases (constructive interference) and some have negative phases (destructive interference). 
A PP machine—a theoretical classical computer that non-deterministically guesses answers and only needs to be right slightly more than half the time—is perfectly engineered for this kind of mathematical balancing act.

We can cleverly program the PP machine so that its accepting and rejecting branches physically mirror the positive and negative phase contributions of the quantum circuit \cite{adleman1997}. 
Because the PP class allows for an incredibly razor-thin margin of error, it can absorb the bounded 66\% success rate of the BQP machine and translate it into a classical probability that simply floats above 50\%. 
In short, the theoretical power of unbounded classical guessing is vast enough to mathematically simulate quantum wave interference.
\end{proof}

\textbf{Theorem 5:} $\text{PP} \subseteq \text{PSPACE}$
\begin{proof}
Finally, can we simulate that incredibly powerful PP guessing machine using limited memory? 
Absolutely, by using the exact same depth-first search trick we applied in Theorem 3. 

The probabilistic PP machine generates a massive decision tree of random coin flips. 
Instead of evaluating the whole tree at once, a deterministic classical machine can methodically walk down one branch at a time. 
It keeps a running classical counter of how many branches ultimately end in an "accept" state. 
Once it painstakingly finishes checking every single branch, it simply asks one question: is the final counter greater than half of the total branches? 

Keeping track of the current path and maintaining a single numerical counter only takes a polynomial amount of memory space. 
Therefore, the PP class fits securely inside the bounds of PSPACE.
\end{proof}